
% !TEX root = ../aa_SoM_Chapters.tex

%% HARRIS: Chapter 10

% ö = \%c3\%b6
% ä = \%C3\%A4

\xdef\sectiontitle{\IIsectiontitle} %aiheuttama

%%%%%%%%%%%%%%%
\begin{frame}[label=2_6_a]%[label=2_5_TOC1]
\frametitle{\texorpdfstring{\culine[\sectiontitleulinecolor]{\sectiontitle}}{\sectiontitle} }

%\vspace{0.5cm}
%\begin{itemize}[leftmargin=0.5cm,itemsep=3mm]
%    \item[2.1] \hyperlink{2_1_a}{\IIaSubsectiontitle}
%    \item[2.2] \hyperlink{2_2_a}{\IIbSubsectiontitle}	
%    \item[2.3] \hyperlink{2_3_a}{\IIcSubsectiontitle}	
%    \item[2.4] \hyperlink{2_4_a}{\IIdSubsectiontitle}
%    \item[2.5] \hyperlink{2_5_a}{\CDAlert[blue]{\IIeSubsectiontitle}}
%    \item[2.6] \hyperlink{2_6_a}{\IIfSubsectiontitle}
%    \item[2.7] \hyperlink{2_7_a}{\IIgSubsectiontitle}
%    \item[2.8] \hyperlink{2_8_a}{\IIhSubsectiontitle}
%    \item[2.9] \hyperlink{2_9_a}{\IIiSubsectiontitle}
%\end{itemize}

\vspace{0.2cm}
\begin{itemize}[leftmargin=0.5cm,itemsep=3mm]
    \item[2.1] \hyperlink{2_1_a}{\IIaaSubsectiontitle}
    \item[2.2] \hyperlink{2_2_a}{\IIaSubsectiontitle}
    \item[2.3] \hyperlink{2_3_a}{\IIbSubsectiontitle}	
    \item[2.4] \hyperlink{2_4_a}{\IIcSubsectiontitle}	
    \item[2.5] \hyperlink{2_5_a}{\IIdSubsectiontitle}
    \item[2.6] \hyperlink{2_6_a}{\CDAlert[blue]{\IIeSubsectiontitle}}
    \item[2.7] \hyperlink{2_7_a}{\IIfSubsectiontitle}
    \item[2.8] \hyperlink{2_8_a}{\IIgSubsectiontitle}
    \item[2.9] \hyperlink{2_9_a}{\IIhSubsectiontitle}
    \item[2.10] \hyperlink{2_10_a}{\IIiSubsectiontitle}
\end{itemize}

\endofslide \end{frame}
%%%%%%%%%%%%%%%

\xdef\sectiontitle{\IIeSubsectiontitle} 

%%%%%%%%%%%%%%%
\begin{frame}
\frametitle{\texorpdfstring{\culine[\sectiontitleulinecolor]{\sectiontitle}}{\sectiontitle} }

\begin{itemize}
	\item When atoms come together\appear{, the electrons strongly bound to their atomic nuclei remain isolated in their atoms, and do not participate in electronic band formation (Figure)}

\item Higher states\appear{, that's where the outmost electrons} \appear{i.e. the valence electron are,} \appear{are however broken into bands by the ionic potential}

\item In a crystal of $N$ atoms\appear{, energy bands are split into $N$ spatial states} 

\item Due to spin\appear{, the total number of electrons in the band is $2 N$}

\end{itemize}

% 6.5 cm = midway, 10 cm = 3/4 
\vspace{2.5mm}
\begin{minipage}{5.5cm}
\begin{itemize}
\item The difference between insulators \appear{and conductors is found in the \CDAlert[red]{occupation} of the \CDAlert[red]{highest energy band}}\appear{, and whether it's full} \appear{or only partially filled with electrons}
\end{itemize}
\end{minipage}

\xdef\loc{south east}
\begin{tikzpicture}[remember picture,overlay] % anchor=south
    \node (A) [yshift=-0.25*\figYshift,xshift=\figXshift, anchor=\loc] at (current page.\loc) {\includegraphics[page=1,width=8cm]{Figures_booklet/SOM_10_Bonding_figures/SOM_Bonding_Figure_29.png}}; %scale=\figscale. % 8cm max height
\end{tikzpicture}

\endofslide 
%\endofsection
\end{frame}
%%%%%%%%%%%%%%%

%%%%%%%%%%%%%%%%
%\begin{frame}
%\frametitle{\texorpdfstring{\culine[\sectiontitleulinecolor]{\sectiontitle}}{\sectiontitle} }
%
%\begin{itemize}
%	\item Let's look at the periodic potential ...
%	\item $\ldots$ and the possible wave functions\appear{, here $N=8$}\appear{, as an example (only 4 functions presented):}
%	
%	\item We obtained energy bands, which may be overlapping with each other (depending on $a$)
%	
%\item And we found out that a great fraction of the energy of electrons is kinetic\appear{, potential energy deforms the curves, and gaps are formed}
%
%\end{itemize}
%
%
%\endofslide 
%%\endofsection
%\end{frame}
%%%%%%%%%%%%%%%%

%%%%%%%%%%%%%%%
\begin{frame}
\frametitle{\texorpdfstring{\culine[\sectiontitleulinecolor]{\sectiontitle}}{\sectiontitle} }
\framesubtitle{Real DOS}

\begin{itemize}
%	\item To study the band structure, we can take approaches of different level in details 
	\item Next, we need to calculate how the electrons in the solid will populate the energy band
	\appear{\xdef\loc{south east}%
\begin{tikzpicture}[remember picture,overlay]% anchor=south
    \node (A) [yshift=-0.15*\figYshift,xshift=\figXshift, anchor=\loc] at (current page.\loc) {\includegraphics[page=1,width=8cm]{Figures_booklet/SOM_10_Bonding_figures/Tikz_SOM_Bonding_EFermi_schematic}};%scale=\figscale. % 8cm max height
\end{tikzpicture}}
	\item Recall, that we have already calculated the \CDAlert[fuchsia]{density of states} for a 3D quantum well:
\[
D(E) \equal \frac{1}{8} \appear{4 \pi n^{2} }\frac{d n}{d E}  \equal \frac{(2 s+1) m^{3 / 2} V}{\pi^{2} \hbar^{3} \sqrt{2}} \appear{\sqrt{E}}
\]
\vspace{2.5mm}
\appear{and the associated \CDAlert[fuchsia]{Fermi energy $E_{F}$} was given by}
\end{itemize}
% 6.5 cm = midway, 10 cm = 3/4 
%\vspace{2.5mm}
\begin{minipage}{6.5cm}
\begin{itemize}
	\item[] \[
 E_{F} \equal \frac{\pi^{2} \hbar^{2}}{m} \appear{\textrm{[[}}\frac{3}{(2 s+1) \pi \sqrt{2}} \frac{N}{V} \appear{\textrm{]]}^{2 / 3}}
\]
\end{itemize}
\end{minipage}

\endofslide 
%\endofsection
\end{frame}
%%%%%%%%%%%%%%%

%%%%%%%%%%%%%%%
\begin{frame}
\frametitle{\texorpdfstring{\culine[\sectiontitleulinecolor]{\sectiontitle}}{\sectiontitle} }

\begin{itemize}
	\item In what follows, we could have 3 different cases for density of states:
\item[] \begin{itemize}
	\item density of states $D$ is kept constant \appear{(for simplicity of presentation, choice by Harris)}
\item $D=D(E)$ is derived for square 3D-potential well structures \appear{(on previous lectures)}
\item $D=D(E)$ is derived for Coulombic well structures \appear{(more sophisticated)}

\end{itemize}

\item Additionally, we study temperature cases where:
\item[] \begin{itemize}
	\item $T=0 \mathrm{~K}$
	\item $T>0 \mathrm{~K}$
	\end{itemize}
	
	\item Thus we will have 6 different approaches\appear{, see Table in the next slide}

\end{itemize}

\endofslide 
%\endofsection
\end{frame}
%%%%%%%%%%%%%%%

%%%%%%%%%%%%%%%
\begin{frame}
\frametitle{\texorpdfstring{\culine[\sectiontitleulinecolor]{\sectiontitle}}{\sectiontitle} }

\begin{itemize}
	
\item Please note, that these approaches are used rather randomly in attempt to illustrate the same point 
\implies{ \Alert[fuchsia]{Energy band structure of electrons in a crystalline solid}}
\vspace{4mm}

{ \small
\begin{tabular}{|l|c|l|} 
\hline \Alert[blue]{Density of states, $D$} & FD-statistics at \Alert[red]{$T=0 \mathrm{~K}$} & FD-statistics at \Alert[tunired]{$T>0 \mathrm{~K}$} \\
\hline Constant & X & X \\
\hline Periodic 3D-square well & X & X \\
\hline Periodic Coulombic potential & & \\
\hline
\end{tabular}
}

\end{itemize}

\xdef\loc{south}%
\begin{tikzpicture}[remember picture,overlay]% anchor=south
    \node (A) [yshift=-0.25*\figYshift,xshift=0*\figXshift, anchor=\loc] at (current page.\loc) {\includegraphics[page=1,width=8cm]{Figures_booklet/SOM_10_Bonding_figures/Tikz_SOM_Bonding_EFermi_schematic}};%scale=\figscale. % 8cm max height
\end{tikzpicture}

\endofslide 
%\endofsection
\end{frame}
%%%%%%%%%%%%%%%

%%%%%%%%%%%%%%%
\begin{frame}
\frametitle{\texorpdfstring{\culine[\sectiontitleulinecolor]{\sectiontitle}}{\sectiontitle} }
\framesubtitle{Band overlap}

\begin{itemize}
	\item Let's then check the energy levels of electrons for light elements in the periodic table:
	% 6.5 cm = midway, 10 cm = 3/4 
\end{itemize}
\vspace{1.5mm}

\begin{minipage}{3.75cm} \small
\begin{itemize}
\item[\ii] \CDAlert[red]{Lithium and Sodium} \appear{both have one electron on the valence states $2 \mathrm{s}$ and $3 \mathrm{s}$, respectively.} \appear{So, the bands are 'half full'.} \appear{They are assumed to be conductors.} \appear{And, in fact, they are \CDAlert[red]{conductors}}

\item[\ii] \CDAlert[blue]{Beryllium and Magnesium} \appear{have two electrons in the states $2 \mathrm{s}$ and $3s$, respectively}. \appear{So, for them, the highest occupied band is 'full'.} \appear{It could be \CDAlert[blue]{assumed} that they are \CDAlert[blue]{insulators}}
\end{itemize}
\end{minipage}


\xdef\loc{south east}
\begin{tikzpicture}[remember picture,overlay] % anchor=south
    \node (A) [yshift=-0.25*\figYshift,xshift=\figXshift, anchor=\loc] at (current page.\loc) {\includegraphics[page=1,width=10cm]{Figures_booklet/SOM_10_Bonding_figures/SOM_Atomic_Figure_14.png}}; %scale=\figscale. % 8cm max height
\end{tikzpicture}

\endofslide 
%\endofsection
\end{frame}
%%%%%%%%%%%%%%%

%%%%%%%%%%%%%%%
\begin{frame}
\frametitle{\texorpdfstring{\culine[\sectiontitleulinecolor]{\sectiontitle}}{\sectiontitle} }
\framesubtitle{Band overlap}

\begin{itemize}
\item Relating atomic states $ \textrm{((}n, \ell, m_{\ell} \textrm{))}$ to bands in a solid can be complicated. \appear{Here, we treat the issue qualitatively, considering at first that there is no overlapping of different bands}

\item The \CDAlert[red]{highest energy band}, that would be \Alert[red]{completely full at zero kelvins}\appear{, is defined as \CDAlert[red]{valence band}}
	% 6.5 cm = midway, 10 cm = 3/4 

\end{itemize}
\vspace{2.5mm}

\begin{minipage}{7.5cm}
\begin{itemize}
\item The \CDAlert[blue]{next} (higher) \Alert[blue]{allowed band is called} \CDAlert[blue]{conduction band}

\item For lithium, as an example\appear{, the $2 \mathrm{s}$ band is a conduction band}
\item For sodium\appear{, the $3 \mathrm{s}$ band is a conduction band}
\end{itemize}
\end{minipage}

%% cropped image
\xdef\loc{south east}
\begin{tikzpicture}[remember picture,overlay] % anchor=south
    \node (A) [yshift=-0.25*\figYshift,xshift=\figXshift, anchor=\loc] at (current page.\loc) {\includegraphics[page=1,width=6cm]{Figures_booklet/SOM_10_Bonding_figures/SOM_Bonding_Figure_30_cropped.png}}; %scale=\figscale. % 8cm max height
\end{tikzpicture}

\endofslide 
%\endofsection
\end{frame}
%%%%%%%%%%%%%%%

%%%%%%%%%%%%%%%
\begin{frame}
\frametitle{\texorpdfstring{\culine[\sectiontitleulinecolor]{\sectiontitle}}{\sectiontitle} }
\framesubtitle{Band overlap}

\begin{itemize}
	\item The information above\appear{, however is still not sufficient to decide whether an element is an insulator or a conductor}

\item One needs also knowledge about the \Alert[red]{actual energies} \appear{at the above-mentioned states}\appear{, or \CDAlert[red]{'energy bands'}}
\end{itemize}

% 6.5 cm = midway, 10 cm = 3/4 
\vspace{2.5mm}
\begin{minipage}{4.5cm}
\begin{itemize}
	\item Fig.~30, erroneously\appear{, assumes that there is an energy band gap between the filled $2 \mathrm{s}$ band} \appear{and the above $2 \mathrm{p}$ for Be} 
	\item If Fig.~30 would be true, \appear{Be would be an insulator (no free electrons)}
\end{itemize}
\end{minipage}


\xdef\loc{south east}
\begin{tikzpicture}[remember picture,overlay] % anchor=south
    \node (A) [yshift=-0.25*\figYshift,xshift=\figXshift, anchor=\loc] at (current page.\loc) {\includegraphics[page=1,width=9cm]{Figures_booklet/SOM_10_Bonding_figures/SOM_Bonding_Figure_30.png}}; %scale=\figscale. % 8cm max height
\end{tikzpicture}

\endofslide 
%\endofsection
\end{frame}
%%%%%%%%%%%%%%%

%%%%%%%%%%%%%%%
\begin{frame}
\frametitle{\texorpdfstring{\culine[\sectiontitleulinecolor]{\sectiontitle}}{\sectiontitle} }
\framesubtitle{Band overlap}

\begin{itemize}
	\item Both \CDAlert[red]{Beryllium} and \CDAlert[red]{Magnesium}, however, \CDAlert[red]{are electrical conductors}

\item The point is, that the \CDAlert[fuchsia]{energies of $2p$ (and $3 \mathrm{p}$)} band for beryllium (and magnesium) are \CDAlert[fuchsia]{overlapping with 2s (3s)} 

\item Therefore, even the electrons in the full 2s band of beryllium \appear{(full 3s band of magnesium)} \appear{have access to}
% 6.5 cm = midway, 10 cm = 3/4 
\end{itemize}
%\vspace{2.5mm}
\begin{minipage}{7.1cm}
\begin{itemize}
	\item[]  $2 p$ levels (3p levels): 
	\implies{ possibility to move around\appear{, and form a \\
		      free electron quantum gas} } 
	\implies{ and to enable \CDAlert[blue]{electrical conductivity} }
\end{itemize}
\end{minipage}


\xdef\loc{south east}
\begin{tikzpicture}[remember picture,overlay] % anchor=south
    \node (A) [yshift=-0.25*\figYshift,xshift=\figXshift, anchor=\loc] at (current page.\loc) {\includegraphics[page=1,width=6.5cm]{Figures_booklet/SOM_10_Bonding_figures/SOM_Bonding_Figure_31.png}}; %scale=\figscale. % 8cm max height
\end{tikzpicture}

\endofslide 
%\endofsection
\end{frame}
%%%%%%%%%%%%%%%

%%%%%%%%%%%%%%%
\begin{frame}
\frametitle{\texorpdfstring{\culine[\sectiontitleulinecolor]{\sectiontitle}}{\sectiontitle} }
\framesubtitle{Memorize Figure!}

\seminormal
\begin{itemize}
	\item Figure 32 illustrates the difference between \Alert[red]{insulators, conductors and semiconductors}\appear{, first at $T=0 \mathrm{~K}$ (left)}\appear{, then at $\mathrm{T}>0 \mathrm{~K}$ (right).} \appear{Also enclosed are the Fermi–Dirac distributions at $T=0 \mathrm{~K}$} \appear{(practically a step function up to $E_{F}$)} \appear{and at $\mathrm{T}>0 \mathrm{~K}$ (smoother transition at $E_{F}$).} \appear{Material is a conductor}\appear{, if the band occupied by the highest-energy electrons is not full}\appear{, making it by definition a conduction band}

\item Even at $T=0~\mathrm{K}$\appear{, the electrons are free to absorb energy from external field}
\end{itemize}

\xdef\loc{south}
\begin{tikzpicture}[remember picture,overlay] % anchor=south
    \node (A) [yshift=-0.25*\figYshift,xshift=\figXshift, anchor=\loc] at (current page.\loc) {\includegraphics[page=1,width=13cm]{Figures_booklet/SOM_10_Bonding_figures/SOM_Bonding_Figure_32.png}}; %scale=\figscale. % 8cm max height
\end{tikzpicture}

\endofslide 
%\endofsection
\end{frame}
%%%%%%%%%%%%%%%

%%%%%%%%%%%%%%%
\begin{frame}
\frametitle{\texorpdfstring{\culine[\sectiontitleulinecolor]{\sectiontitle}}{\sectiontitle} }
\framesubtitle{ Insulators}

\begin{itemize}
	\item If a solid has a completely filled valence band\appear{, and the gap between the valence band} \appear{and empty conduction band} \appear{is larger than about $2 \mathrm{~eV}$}\appear{, the solid is an insulator} 
	\item The next-higher allowed energy band \appear{(conduction band)} \appear{is so high} \appear{that it cannot be reached by ordinary external fields}\appear{, and the electrons cannot occupy any states they were not occupying already} 
	\implies{ no response to the external field }

\item E.g. \CDAlert[blue]{ionic crystals are insulators} 
\item The energy bands of $\mathrm{NaCl}$ arise from levels of the Na$^+$ and Cl$^-$ ions. \appear{Both ions have closed-shell configuration}\appear{, the highest occupied band is full, while the next band up is empty}

\end{itemize}

\endofslide 
%\endofsection
\end{frame}
%%%%%%%%%%%%%%%

%%%%%%%%%%%%%%%
\begin{frame}
\frametitle{\texorpdfstring{\culine[\sectiontitleulinecolor]{\sectiontitle}}{\sectiontitle} }
\framesubtitle{Group 14}

\semismall

\begin{itemize}
	 \item Carbon has two 2s and $2p$ electrons. \appear{In principle, the p-level is unfilled.} \appear{However,} %\appear{However, the $2 \mathrm{s}$ and $2 \mathrm{p}$ states hybridize} \appear{(into sp$^3$)} \appear{when carbon forms the covalent diamond structure (Figure)}
\end{itemize}

% 6.5 cm = midway, 10 cm = 3/4 
%\vspace{2.5mm}
\begin{minipage}{7.4cm}
\begin{itemize}
	\item[] the \CDAlert[fuchsia]{$2 \mathrm{s}$ and $2 \mathrm{p}$ states hybridize} \appear{(into sp$^3$)} \appear{when carbon forms the covalent diamond structure (Figure)}
	
	%\item In figure, the s- and p-levels split \appear{due to small atomic separation}\appear{, giving rise to hybridized sp$^3$ orbitals} \appear{and formation of covalent diamond structure.} %\appear{The energy of four space symmetric states} \appear{(one for $2 \mathrm{s}$, three for $2 \mathrm{p}$)} \appear{is lowered}\appear{, and the energy of four states is raised}

\item The diamond lattice spacing is $0.154 \mathrm{~nm}$. \appear{The energy gap between filled valence band and empty conduction band is $\approx 5.4 \mathrm{~eV}$.} \appear{At temperature of $300 \mathrm{~K}$}\appear{, only few electrons can reach the conduction band }
\implies{The valence band contains four electrons\\ 
		(per atom)\appear{, while the conduction band\\ 
		is empty}}
\implies{\CDAlert[red]{Diamond is an insulator}}	


\item For \CDAlert[blue]{Si and Ge}\appear{, the band gaps are $1.1 \mathrm{~eV}$ and $0.7 \mathrm{~eV}$, respectively}\appear{, making them \CDAlert[blue]{(intrinsic)} \CDAlert[blue]{semiconductors}}
\end{itemize}
\end{minipage}
%\vspace{2.5mm}



\xdef\loc{south east}
\begin{tikzpicture}[remember picture,overlay] % anchor=south
    \node (A) [yshift=-0.25*\figYshift,xshift=\figXshift, anchor=\loc] at (current page.\loc) {\includegraphics[page=1,width=6.3cm]{Figures_booklet/SOM_Tipler_Solid_state_physics_figures/SoM_Bonding_of_solids_Tipler_Band_hybridization_C_Si_Ge.png}}; %scale=\figscale. % 8cm max height
\end{tikzpicture}


\endofslide 
%\endofsection
\end{frame}
%%%%%%%%%%%%%%%

%%%%%%%%%%%%%%%
\begin{frame}
\frametitle{\texorpdfstring{\culine[\sectiontitleulinecolor]{\sectiontitle}}{\sectiontitle} }
\framesubtitle{Semiconductors}

\begin{itemize}

	\item If the \CDAlert[fuchsia]{gap} between the valence band and empty conduction band \CDAlert[fuchsia]{is small, the material is a semiconductor}

\item Note that semiconductors are insulators at $T=0 \mathrm{~K}$\appear{, but some electrons are excited to the conduction band at normal (room) temperatures}\appear{, leaving holes in the valence band }

\item However, the number of these electrons is small \appear{(compared to conductors)}

\item The Fermi level $E_F$ is approximately in the middle of the gap

\item These materials are \CDAlert[fuchsia]{intrinsic semiconductors}\appear{, intrinsic meaning} \appear{that they do not contain dopants} \appear{or impurities} \appear{(more information on semiconductors follows also in the later sections)}
\end{itemize}

\endofslide 
%\endofsection
\end{frame}
%%%%%%%%%%%%%%%

%%%%%%%%%%%%%%%%
%\begin{frame}
%\frametitle{\texorpdfstring{\culine[\sectiontitleulinecolor]{\sectiontitle}}{\sectiontitle} }
%
%\begin{itemize}
%	\item More detailed presentation of sodium's energy bands
%\end{itemize}
%
%%\xdef\loc{east}
%%\begin{tikzpicture}[remember picture,overlay] % anchor=south
%%    \node (A) [yshift=\figYshift,xshift=\figXshift, anchor=\loc] at (current page.\loc) {\includegraphics[page=1,height=7cm]{Figures_booklet/SOM_10_Bonding_figures/SOM_Bonding_Figure_1.png}}; %scale=\figscale. % 8cm max height
%%\end{tikzpicture}
%
%\endofslide 
%%\endofsection
%\end{frame}
%%%%%%%%%%%%%%%%



%%%%%%%%%%%%%%%
\begin{frame}
\frametitle{\texorpdfstring{\culine[\sectiontitleulinecolor]{\sectiontitle}}{\sectiontitle} }
\framesubtitle{Fermi level}
%The Conductivity Gap

\begin{itemize}
	\item Electrons obey Fermi–Dirac statistics\appear{, and their occupation number is distributed to different energies according to:}
\[
 \mathbb{N} (E)_{F D}  \equal \frac{1}{e^{ \textrm{((}E-E_{F} \textrm{))} / k_{B} T}+1} \appear{< 1} \qquad \appear{\text{(exclusion principle)}}
\]
\end{itemize}
\vspace{2.5mm}
\begin{minipage}{7cm}
\begin{itemize}
%which is always less than 1 (exclusion principle). 
	\item At higher $T$, some electrons are excited to higher energies 
	
	\item The distribution remains symmetric around $E_{F}$ \appear{and so the number of electrons appearing in the conduction band needs to be equal to the number of electrons disappearing from the valence band} 
	\implies{ \Alert[fuchsia]{Fermi energy $E_F$ is in the middle \\
			of the gap} }
\end{itemize}
\end{minipage}


\xdef\loc{south east}
\begin{tikzpicture}[remember picture,overlay] % anchor=south
    \node (A) [yshift=-0.25*\figYshift,xshift=\figXshift, anchor=\loc] at (current page.\loc) {\includegraphics[page=1,width=6.5cm]{Figures_booklet/SOM_10_Bonding_figures/SOM_Bonding_Figure_33.png}}; %scale=\figscale. % 8cm max height
\end{tikzpicture}

\endofslide 
%\endofsection
\end{frame}
%%%%%%%%%%%%%%%

%%%%%%%%%%%%%%%
\begin{frame}
\frametitle{\texorpdfstring{\culine[\sectiontitleulinecolor]{\sectiontitle}}{\sectiontitle} }
\framesubtitle{Band gap}

\begin{itemize}
	\item Note, here in figures\appear{, the density of states $D$ is assumed to be constant} \appear{(simplification by Harris)}
	% 6.5 cm = midway, 10 cm = 3/4 
\end{itemize}
\vspace{2.5mm}

\begin{minipage}{7.3cm}
\begin{itemize}
	%\item[] assumed to be a constant \appear{(presentation by Harris)} 
	
	\item Using $D=$~constant, we quickly obtain the number of electrons per energy range $d E$\appear{, that is $ \mathbb{N} (E)_{F D} \times D$} \appear{i.\,e. "the number of electrons per state times the number of states per $d E$}
	
	\item Once we know the number of electrons per energy range\appear{, we can calculate how the electron band structure is occupied by electrons}
	\implies{ Once level occupation is known\appear{, one can\\ 
		      predict whether the solid is an insulator}\appear{, \\
		      a semiconductor}\appear{ or a metal} }	
\end{itemize}
\end{minipage}

\xdef\loc{south east}
\begin{tikzpicture}[remember picture,overlay] % anchor=south
    \node (A) [yshift=-0.25*\figYshift,xshift=\figXshift, anchor=\loc] at (current page.\loc) {\includegraphics[page=1,height=7.3cm]{Figures_booklet/SOM_10_Bonding_figures/SOM_Bonding_Figure_34.png}}; %scale=\figscale. % 8cm max height
\end{tikzpicture}

\endofslide 
%\endofsection
\end{frame}
%%%%%%%%%%%%%%%

%%%%%%%%%%%%%%%
\begin{frame}
\frametitle{\texorpdfstring{\culine[\sectiontitleulinecolor]{\sectiontitle}}{\sectiontitle} }
\framesubtitle{Band occupation}

\begin{itemize}
	\item Now we estimate the number of electrons which are excited over the gap \appear{($D=$~constant)}

\item First we can calculate the number of \CDAlert[red]{electrons in the valence band} at $T=0$:
\[
N_{V}  \equal \appear{\nint_{E_{V, bottom}}^{E_{V, top}}}  \appear{\mathbb{N} (E)_{F D}} \appear{D} \appear{d E} \equal \appear{\nint_{E_{V, bottom}}^{E_{V, top}}} \appear{1 \cdot D} \appear{d E} \equal \appear{D \Delta E_{V}}
\]

\item Then we find the \CDAlert[blue]{number of excited electrons} at $T>0$:
\[
N_{\text {excited }}  \equal \appear{\nint_{E_{F}+\frac{1}{2} E_{g a p}}^{\infty}} \frac{1}{e^{ \textrm{((}E-E_{F} \textrm{))} / k_{B} T}+1} \appear{D} \appear{d E} \equal \appear{D k_{B} T} \appear{\ln  \textrm{[[}}\appear{1+e^{-E_{g a p} / 2 k_{B} T}} \appear{\textrm{]]}}
\]
\appear{recalling that} $\appear{\ln (1+\varepsilon)} \approx \appear{\varepsilon}\appear{\,, \text { if } \varepsilon \ll 1} $\appear{, we get}
$$
\Rightarrow \qquad \appear{N_{\text {excited}}} \approx \appear{D} \appear{k_{B}} \appear{T} \appear{e^{-E_{gap} / 2 k_{B} T}}
$$
\end{itemize}

\endofslide 
%\endofsection
\end{frame}
%%%%%%%%%%%%%%%

%%%%%%%%%%%%%%%
\begin{frame}
\frametitle{\texorpdfstring{\culine[\sectiontitleulinecolor]{\sectiontitle}}{\sectiontitle} }
\framesubtitle{Band occupation}

\begin{itemize}
	\item Thus, the ratio between electrons excited at $T>0$ \appear{and the total number in the valence band at at $T=0$}\appear{, will be:}
\[
\frac{N_{\text {excited }}}{N_{V}}  \equal  \frac{k_{B} T}{\Delta E_{V}} \appear{e^{-E_{g a p} / 2 k_{B} T}}
\]
\item The energy bands tend to be usually ${\sim}10$~eV high\appear{, and $k_{B} T \cong \frac{1}{40} \mathrm{~eV}$ et room $T$.} \appear{So the multiplicative factor in the equation above is roughly $10^{-3}$}

\item But, in fact the \CDAlert[fuchsia]{exponential factor dominates}:
\item[] \begin{itemize}
	\item If the band gap is $5.4 \mathrm{~eV}$ (carbon)\appear{, then we obtain the exponential factor to be $10^{-42}$}
\item If the band gap is $1.1 \mathrm{~eV}$ (silicon)\appear{, then we obtain the exponential factor to be $10^{-8}$}
\end{itemize}

\end{itemize}

%\xdef\loc{east}
%\begin{tikzpicture}[remember picture,overlay] % anchor=south
%    \node (A) [yshift=\figYshift,xshift=\figXshift, anchor=\loc] at (current page.\loc) {\includegraphics[page=1,height=7cm]{Figures_booklet/SOM_10_Bonding_figures/SOM_Bonding_Figure_1.png}}; %scale=\figscale. % 8cm max height
%\end{tikzpicture}

\endofslide 
%\endofsection
\end{frame}
%%%%%%%%%%%%%%%

%%%%%%%%%%%%%%%
\begin{frame}
\frametitle{\texorpdfstring{\culine[\sectiontitleulinecolor]{\sectiontitle}}{\sectiontitle} }
\framesubtitle{Example}

\seminormal
%Example 6.
\begin{itemize}
	\item \textbf{Q}: \appear{Near room-$T$}\appear{, an increase of $1 $~K decreases the conductivity of copper by $\approx 0.4 \%$.} \appear{For an unknown semiconductor}\appear{, $\Delta T=+1 K$ increases the conductivity by $4.8 \%$.} \appear{Estimate the semiconductor's band gap.}
	
\item \textbf{A}: \appear{Let's estimate the number of electrons at excited states both at $300 \mathrm{~K}$} \appear{and $301 \mathrm{~K}$}\appear{, and then set the ratio to $1.048$}
\[ \semismall
\frac{N_{\text{exc,i}}}{N_{V}} \equal \frac{k_{B} T_{i}}{\Delta E_{V}} \appear{e^{-E_{g a p} / 2 k_{B} T_{i}}} \appear{\,,} \qquad  \qquad  
\frac{N_{\text{exc,f}}}{N_{V}}  \equal \frac{k_{B} T_{f}}{\Delta E_{V}} \appear{e^{-E_{g a p} / 2 k_{B} T_{f}}}
\]

$$ \semismall
\begin{aligned}
&\Rightarrow& \qquad \frac{ N_{\text{exc,f}} \bg[1.3]{/} \,N_{V}}{ N_{\text {exc,i}} /\, N_{V}}& \equal \frac{T_{f}}{T_{i}} \appear{e^{- \textrm{[[}}\frac{E_{gap}}{2 k_{B} T_{f}} \appear{\textrm{]]}+ \textrm{[[}}\frac{E_{g a p}}{2 k_{B} T_{i}}\appear{ \textrm{]]}}} \equal \appear{1.048} \\
&\Rightarrow& \qquad \qquad \frac{E_{gap}}{2 k_{B}} \appear{\textrm{[[}}\frac{T_{f}-T_{i}}{T_{i} T_{f}} \appear{\textrm{]]}} & \equal \appear{\ln  \textrm{[[}}\appear{1.048} \frac{T_{i}}{T_{f}} \appear{\textrm{]]}} \\
 &\Rightarrow& \qquad  \appear{E_{gap}}& \equal \appear{2 k_{B}} \appear{\textrm{[[}}\frac{T_{i} T_{f}}{T_{f}-T_{i}} \appear{\textrm{]]}} \appear{\ln  \textrm{[[}1.048} \frac{T_{i}}{T_{f}}\appear{ \textrm{]]}}  \equal \appear{0.68 ~eV} \qquad 
\end{aligned}
$$
\appear{Obtained band gap $(\sim 0.7 \mathrm{eV}$) matches closely to the $E_{gap}$ of \CDAlert[blue]{germanium}}
\end{itemize}

%\endofslide 
\endofsection
\end{frame}
%%%%%%%%%%%%%%%
